\documentclass[tikz,border=5pt]{standalone}
\usepackage[T1]{fontenc}
\usepackage[utf8]{inputenc}
\usepackage[italian]{babel}
\usepackage{amsmath}
\usetikzlibrary{arrows.meta,calc,decorations.pathreplacing,positioning}

\begin{document}
\begin{tikzpicture}[
    % Figura generale ispirata alla slide 16.
    cell/.style={
        draw,
        minimum width=0.38cm,
        minimum height=0.50cm,
        font=\ttfamily\small,
        inner sep=0pt
    },
    search/.style={fill=blue!7},
    lookahead/.style={fill=green!9},
    match/.style={fill=orange!28},
    pointer/.style={-Latex,thick},
    note/.style={font=\small,align=center},
    title/.style={font=\bfseries\small}
]

% Riga dei caratteri: abbaccbbcabc | bbcacbcaacbbcb.
\foreach \ch [count=\i from 1] in {
    a,b,b,a,c,c,b,b,c,a,b,c,
    b,b,c,a,c,b,c,a,a,c,b,b,c,b
} {
    \ifnum\i<13
        \node[cell,search] (c\i) at ({0.38*(\i-1)},0) {\ch};
    \else
        \node[cell,lookahead] (c\i) at ({0.38*(\i-1)},0) {\ch};
    \fi
}

% Match bbca nel search buffer e all'inizio del look-ahead buffer.
\foreach \i/\ch in {7/b,8/b,9/c,10/a,13/b,14/b,15/c,16/a} {
    \node[cell,match] (c\i) at (c\i) {\ch};
}

% Graffe descrittive.
\draw[decorate,decoration={brace,mirror,amplitude=4pt}]
    ($(c1.south west)+(0,-0.09)$) -- ($(c12.south east)+(0,-0.09)$)
    node[midway,below=5pt,title,blue!60!black] {search buffer};
\draw[decorate,decoration={brace,mirror,amplitude=4pt}]
    ($(c13.south west)+(0,-0.09)$) -- ($(c26.south east)+(0,-0.09)$)
    node[midway,below=5pt,title,green!45!black] {look-ahead buffer};

% Puntatori i e j.
\draw[pointer,blue!70!black] ($(c13.north)+(0,0.66)$) -- ($(c13.north)+(0,0.10)$)
    node[pos=0,above,note] {$i$};
\draw[pointer,orange!75!black] ($(c7.north)+(0,0.66)$) -- ($(c7.north)+(0,0.10)$)
    node[pos=0,above,note] {$j$};

\end{tikzpicture}
\end{document}
